\documentclass{article}

% Language setting
\usepackage[english]{babel}

% Set page size and margins
\usepackage[letterpaper,top=2cm,bottom=2cm,left=3cm,right=3cm,marginparwidth=1.75cm]{geometry}

% Useful packages
\usepackage[utf8]{inputenc}
\usepackage{latexsym}
\usepackage{amsfonts}
\usepackage{amssymb}
\usepackage{amsthm}
\usepackage{amsmath}
\usepackage{tikz}
\usepackage{circuitikz} 
\usepackage{listings}
\usepackage{mathtools}
\usepackage{float}
\usepackage{natbib}
\usepackage{fourier}
\bibpunct{[}{]}{,}{n}{}{,~}

\title{Software for Data Science: Julia}
\date{17/11/2025}

\begin{document}
\maketitle
\section*{Exercise 1}

Assign different values to new variables (\texttt{x}, \texttt{y}, and \texttt{z}), such as an integer, a floating-point number, and a string. What type does Julia infer for each? Use the \texttt{typeof} function to check.

\noindent
Try reassigning one variable (e.g., \texttt{x}) to a value of a different type. What happens to its inferred type?

\section*{Exercise 2}

Declare a new variable with an explicit type using the \texttt{::} operator. For example, force \texttt{x} to be \texttt{Float64}. 

\noindent
Experiment by assigning a value incompatible with the declared type. What error do you see?

\section*{Exercise 3}

Create a tuple that mixes different types (e.g., an integer, a float, and a string). Use the \texttt{typeof} function to inspect the type of the whole tuple.

\noindent
Access the individual elements and use \texttt{typeof} to check their types.

\section*{Exercise 4}

Define a struct called \texttt{Book} with the following fields: title (\texttt{String}), author (\texttt{String}), and pages (\texttt{Int}). Create two instances of \texttt{Book} and print out their fields.

\noindent
Now, define a \texttt{mutable struct} called \texttt{LibraryBook} with the same fields as \texttt{Book} but add a new Boolean field called \texttt{checked\_out} (set to false by default). Create a \texttt{LibraryBook} instance for one of your book titles, change its \texttt{checked\_out} status to true, and print the updated struct.

\noindent
Write a function that takes a \texttt{LibraryBook} as input and prints a message that tells whether the book is checked out or available.


\section*{Exercise 5}

In this exercise the goal is to do a manual conversion from decimal numbers to binary.

\noindent
Pick the following decimal numbers: $13$, $25$, and $42$.

\noindent
For each number, follow these steps by hand:

\begin{itemize}
    \item Divide the number by $2$ and write down the remainder.

    \item Take the integer quotient from the previous division and repeat the process (divide by $2$, write the remainder).

    \item Continue until the quotient is $0$.

    \item The binary representation is the sequence of remainders, read from last to first (bottom to top).
\end{itemize}

\noindent
Example (provided for $13$):

\begin{itemize}
    \item $13 \div 2 = 6$ remainder $1$

    \item $6 \div 2 = 3$ remainder $0$

    \item $3 \div 2 = 1$ remainder $1$

    \item $1 \div 2 = 0$ remainder $1$
    
    \item Binary: Write remainders from bottom to top: $1101$
\end{itemize}

\noindent
Now, complete the steps for $25$ and $42$.

\end{document}